\chapter*{TỔNG KẾT}
\addcontentsline{toc}{chapter}{TỔNG KẾT}

Đề tài đã xây dựng và đánh giá hệ Neuro-Fuzzy thích nghi ANFIS cho bài toán dự báo phụ tải điện hộ gia đình theo giờ. Bài toán được thiết kế cho hai mốc dự báo gồm 1 giờ tới và 24 giờ tới, tương ứng với hai nhãn \texttt{target\_1h} và \texttt{target\_24h}. Dữ liệu đầu vào được xây dựng từ dữ liệu thời tiết theo giờ tại Hà Nội, các hồ sơ sinh hoạt hộ gia đình và chuỗi phụ tải điện mô phỏng.

Trong bản hoàn thiện cuối cùng, mô hình ANFIS sử dụng Core Features gồm \texttt{apparent\_temperature}, \texttt{humidity}, \texttt{hour\_sin}, \texttt{hour\_cos}, \texttt{occupancy\_level}, \texttt{load\_lag\_1} và \texttt{load\_lag\_24}. Việc bổ sung \texttt{load\_lag\_1} giúp mô hình khai thác thêm thông tin phụ tải rất gần thời điểm dự báo. Sau khi huấn luyện lại, ANFIS đạt RMSE 0.187050 và $R^2$ 0.898323 cho horizon \texttt{1h}; với horizon \texttt{24h}, mô hình đạt RMSE 0.231865 và $R^2$ 0.843810.

Kết quả cho thấy ANFIS đạt chất lượng dự báo cạnh tranh và phù hợp với mục tiêu của học phần Tính toán mềm nhờ khả năng kết hợp học từ dữ liệu với diễn giải bằng luật mờ. Tuy nhiên, ANFIS không phải mô hình có RMSE tốt nhất trong toàn bộ thực nghiệm. Các mô hình ensemble như \texttt{Random Forest} và \texttt{XGBoost} đạt hiệu năng tốt hơn khi được huấn luyện trên Extended Dataset, do có khả năng khai thác nhiều đặc trưng mở rộng và các tương tác phi tuyến phức tạp hơn.

Hạn chế chính của đề tài là dữ liệu phụ tải điện vẫn là dữ liệu mô phỏng dựa trên profile hộ gia đình, chưa phải dữ liệu công tơ điện thực tế. Vì vậy, kết quả hiện tại phù hợp để kiểm chứng quy trình và so sánh mô hình trong phạm vi nghiên cứu, nhưng chưa đủ để kết luận trực tiếp về hiệu quả vận hành trong môi trường thực tế.

Trong hướng phát triển tiếp theo, hệ thống nên được đánh giá trên dữ liệu EVN hoặc dữ liệu smart meter thực tế, tối ưu tham số membership functions, tuning số luật mờ, huấn luyện riêng theo từng profile hộ gia đình và thử thêm các mô hình chuỗi thời gian như LSTM hoặc GRU. Các hướng này có thể giúp tăng độ tin cậy thực nghiệm và mở rộng khả năng ứng dụng của hệ thống dự báo phụ tải điện hộ gia đình.
